\begin{table}[h]\small
\caption{La misma metodología con distinto cómputo: CPU propia, GPU de Colab, GPU de Kaggle y artículo}
\label{tab:cpugpu}
\centering
\resizebox{\textwidth}{!}{\begin{tabular}{lcccc}
\toprule
 & CPU propia (8 núcleos) & Google Colab (T4) & Kaggle (T4) & Artículo (GPU)\\
\midrule
Modelo & 128/8, $k=9$, $p=2$ & 256/8, $k=5$, $p=2$ & 256/8, $k=5$, $p=2$ & 256/8, $k=9$ o $5$, $p=2$\\
Parámetros & 223\,370 & 594\,186 & 594\,186 & 709\,k / 594\,k\\
Aumentación & escalado, recorte y espejo & + RandAugment, ColorJitter, borrado & Colab + mixup y CutMix & \texttt{timm} + mixup y CutMix\\
Weight decay & 0,01 & 0,005 & 0 & 0 (Tabla 4)\\
Épocas & \epocasCNN & \epColab & \epKaggle & 200\\
Tiempo por época & \minEpCPU{} min & \segColab{} s & \segKaggle{} s & no informado\\
Tiempo total & \minCMnueve{} min & \minColab{} min & \minKaggle{} min & no informado\\
Exactitud de test & \pct{\accCMnueve} & \pct{\accColab} & \pct{\accKaggle} & \pct{95,0} / \pct{95,21}\\
\bottomrule
\end{tabular}}
\nota{Una corrida por configuración, sin repetir semillas. Colab y Kaggle: se informa la última época, sin elegir la mejor; ambas con una GPU Tesla T4 (en Kaggle la sesión tenía dos, pero el script usa una). Colab entrenó \colabNota. Valores del artículo: Trockman y Kolter (2023, Tabla 4). Antes de la corrida de Colab que se informa hubo otra idéntica que se cortó en la época 84 (94,50 \%); el registro está en \texttt{results/}.}
\end{table}

La Tabla~\ref{tab:cpugpu} resume la comparación que me parece más interesante para discutir. Hay tres lecturas.
Primero, en un tiempo de reloj parecido (entre 1,5 y 2 horas), la GPU entrenó un modelo 2,7 veces más grande, con
muchas más épocas y con la aumentación completa, y ganó unos 12 puntos con el mismo código. Esto relativiza las
comparaciones entre arquitecturas: una parte de lo que se atribuye a ``un modelo mejor'' corresponde a ``un modelo que
recibió más cómputo y una receta más trabajada''. Es justamente la crítica que se le puede hacer al artículo cuando
compara contra números de otros trabajos.

Segundo, la receta cuenta. Pasar de la receta del repositorio (100 épocas, \pct{\accColab}) a la del artículo (200
épocas, mixup, CutMix y sin \emph{weight decay}) subió la exactitud de \pct{\accColab} a \pct{\accKaggle}, es decir, un
punto, y quedó 0,27 puntos por encima del \pct{95,21} publicado. Con una sola semilla por configuración, esa
diferencia con el artículo está dentro de lo que cabe esperar por azar (del orden de 0,2 a 0,3 puntos), así que lo
prudente es decir que la réplica \emph{iguala} el resultado publicado y no que lo supera. Tampoco puedo separar
cuánto aportó cada cambio (las épocas, mixup y CutMix, o el \emph{weight decay}), porque cambié todo junto; sólo el
artículo muestra que quitar mixup y CutMix casi no cambia el resultado (Tabla 3), lo que sugiere que el grueso viene de
las épocas.

Tercero, el costo. Con la misma GPU, la corrida de Kaggle tardó unos \minKaggle{} minutos (\segKaggle{} s por época)
y la de Colab unos \minColab{} minutos con la mitad de las épocas (\segColab{} s por época): en Kaggle cada época
fue más del doble de rápida. No sé a qué atribuirlo con certeza: la GPU es la misma (Tesla T4), pero cambió la
cantidad de procesos que cargan y aumentan las imágenes (4 en Kaggle y 2 en Colab) y, con RandAugment, esa parte corre
en la CPU y puede ser el cuello de botella. En la práctica, Colab gratuito nos negó la GPU por límite de uso a la
mañana siguiente y había cortado una corrida en la época 84; Kaggle ofrece unas 30 horas semanales y permite dejar la
corrida en segundo plano sin tener la pestaña abierta, lo que la hizo más cómoda y confiable para un entrenamiento
largo.
